
\documentclass{article}
%%%%%%%%%%%%%%%%%%%%%%%%%%%%%%%%%%%%%%%%%%%%%%%%%%%%%%%%%%%%%%%%%%%%%%%%%%%%%%%%%%%%%%%%%%%%%%%%%%%%%%%%%%%%%%%%%%%%%%%%%%%%%%%%%%%%%%%%%%%%%%%%%%%%%%%%%%%%%%%%%%%%%%%%%%%%%%%%%%%%%%%%%%%%%%%%%%%%%%%%%%%%%%%%%%%%%%%%%%%%%%%%%%%%%%%%%%%%%%%%%%%%%%%%%%%%
%TCIDATA{OutputFilter=LATEX.DLL}
%TCIDATA{Version=5.50.0.2960}
%TCIDATA{<META NAME="SaveForMode" CONTENT="1">}
%TCIDATA{BibliographyScheme=Manual}
%TCIDATA{Created=Tuesday, September 29, 2026 11:13:03}
%TCIDATA{LastRevised=Tuesday, September 29, 2026 11:13:55}
%TCIDATA{<META NAME="GraphicsSave" CONTENT="32">}
%TCIDATA{<META NAME="DocumentShell" CONTENT="Scientific Notebook\Blank Document">}
%TCIDATA{CSTFile=Math with theorems suppressed.cst}
%TCIDATA{PageSetup=72,72,72,72,0}
%TCIDATA{AllPages=
%F=36,\PARA{038<p type="texpara" tag="Body Text" >\hfill \thepage}
%}


\newtheorem{theorem}{Theorem}
\newtheorem{acknowledgement}[theorem]{Acknowledgement}
\newtheorem{algorithm}[theorem]{Algorithm}
\newtheorem{axiom}[theorem]{Axiom}
\newtheorem{case}[theorem]{Case}
\newtheorem{claim}[theorem]{Claim}
\newtheorem{conclusion}[theorem]{Conclusion}
\newtheorem{condition}[theorem]{Condition}
\newtheorem{conjecture}[theorem]{Conjecture}
\newtheorem{corollary}[theorem]{Corollary}
\newtheorem{criterion}[theorem]{Criterion}
\newtheorem{definition}[theorem]{Definition}
\newtheorem{example}[theorem]{Example}
\newtheorem{exercise}[theorem]{Exercise}
\newtheorem{lemma}[theorem]{Lemma}
\newtheorem{notation}[theorem]{Notation}
\newtheorem{problem}[theorem]{Problem}
\newtheorem{proposition}[theorem]{Proposition}
\newtheorem{remark}[theorem]{Remark}
\newtheorem{solution}[theorem]{Solution}
\newtheorem{summary}[theorem]{Summary}
\newenvironment{proof}[1][Proof]{\noindent\textbf{#1.} }{\ \rule{0.5em}{0.5em}}
\input{tcilatex}

\begin{document}


\FRAME{dtbpFX}{4.9018in}{2.29in}{0pt}{}{}{Plot}{\special{language
"Scientific Word";type "MAPLEPLOT";width 4.9018in;height 2.29in;depth
0pt;display "USEDEF";plot_snapshots TRUE;mustRecompute FALSE;lastEngine
"MuPAD";xmin "-15";xmax "15";xviewmin "-8";xviewmax "8";yviewmin
"-2.5";yviewmax "2.5";viewset"XY";rangeset"X";plottype 4;labeloverrides
1;x-label "t";axesFont "Times New
Roman,12,0000000000,useDefault,normal";numpoints 100;plotstyle
"patch";axesstyle "normal";axestips FALSE;xis \TEXUX{x};var1name
\TEXUX{$x$};function \TEXUX{$\cos \left( 10x\right) $};linecolor
"maroon";linestyle 1;pointstyle "point";linethickness 3;lineAttributes
"Solid";var1range "-15,15";num-x-gridlines 900;curveColor
"[flat::RGB:0x00800000]";curveStyle "Line";rangeset"X";function \TEXUX{$\cos
\left( 11.x\right) $};linecolor "green";linestyle 1;pointstyle
"point";linethickness 3;lineAttributes "Solid";var1range
"-10,10";num-x-gridlines 900;curveColor "[flat::RGB:0x00008000]";curveStyle
"Line";rangeset"X";VCamFile 'TM4YJ70B.xvz';valid_file "T";tempfilename
'TM4YJ701.wmf';tempfile-properties "XPR";}}Notice that there are places
where the waves are in phase, and places where they are not. The
superposition looks like this\FRAME{dtbpFX}{5.0375in}{2.3531in}{0pt}{}{}{Plot%
}{\special{language "Scientific Word";type "MAPLEPLOT";width 5.0375in;height
2.3531in;depth 0pt;display "USEDEF";plot_snapshots TRUE;mustRecompute
FALSE;lastEngine "MuPAD";xmin "-10";xmax "10";xviewmin "-8";xviewmax
"8";yviewmin "-5";yviewmax "5";viewset"XY";rangeset"X";plottype
4;labeloverrides 1;x-label "t";axesFont "Times New
Roman,12,0000000000,useDefault,normal";numpoints 100;plotstyle
"patch";axesstyle "normal";axestips FALSE;xis \TEXUX{x};var1name
\TEXUX{$x$};function \TEXUX{$\cos \left( 10.0x\right) +\cos \left(
11.x\right) $};linecolor "cyan";linestyle 1;pointstyle "point";linethickness
3;lineAttributes "Solid";var1range "-10,10";num-x-gridlines 900;curveColor
"[flat::RGB:0x00008080]";curveStyle "Line";rangeset"X";VCamFile
'TM4YJ70A.xvz';valid_file "T";tempfilename
'TM4YJ702.wmf';tempfile-properties "XPR";}}

Where there is constructive interference, the resulting wave amplitude is
large, where there is destructive interference, the resulting amplitude is
zero. We get a traveling wave who's amplitude varies. We can find the
amplitude function algebraically.

We can write out the entire resultant wave in our usual way. Our two waves
are%
\[
y_{1}=y_{\max }\sin \left( kx-\omega _{1}t+\phi _{o}\right) 
\]%
\[
y_{2}=y_{\max }\sin \left( kx-\omega _{2}t+\phi _{o}\right) 
\]

and the resultant 
\begin{eqnarray*}
y_{r} &=&2y_{\max }\cos \left( \frac{\left( kx-\omega _{2}t+\phi _{o}\right)
-\left( kx-\omega _{1}t+\phi _{o}\right) }{2}\right) \sin \left( \frac{%
kx-\omega _{2}t+\phi _{o}+kx-\omega _{1}t+\phi _{o}}{2}\right) \\
&=&2y_{\max }\cos \left( \frac{kx-2\pi f_{2}t-\left( kx-2\pi f_{1}t\right) }{%
2}\right) \sin \left( \frac{kx-2\pi f_{2}t+kx-2\pi f_{1}t+2\phi _{o}}{2}%
\right) \\
&=&2y_{\max }\cos \left( 2\pi \frac{f_{1}-f_{2}}{2}t\right) \sin \left(
kx-2\pi \frac{f_{1}+f_{2}}{2}t+\phi _{o}\right) \\
&=&\left[ 2y_{\max }\cos \left( 2\pi \frac{f_{1}-f_{2}}{2}t\right) \right]
\sin \left( kx-2\pi \frac{f_{1}+f_{2}}{2}t+\phi _{o}\right)
\end{eqnarray*}

The sine part is a wave, it is a function of position and time. We see that
it has a frequency that is the average of $f_{1}$ and $f_{2}.$ This is the
frequency we hear. But we have another complicated amplitude term, and this
time it is a function of time just as we suspected. The amplitude has its
own frequency that is half the difference of $f_{1}$ and $f_{2}.$ 
\[
A_{\text{resultant}}=2A\cos \left( 2\pi \frac{f_{1}-f_{2}}{2}t\right)
\]%
So the sound amplitude will vary in time for a given position in the medium.

The situation is odder still. We have a cosine function, but it is really an
envelope for the higher frequency motion of the air particles. The air
molecules move back and forth for both the crest and the trough of the
envelope function. \FRAME{dtbpFX}{5.0375in}{2.3531in}{0pt}{}{}{Plot}{\special%
{language "Scientific Word";type "MAPLEPLOT";width 5.0375in;height
2.3531in;depth 0pt;display "USEDEF";plot_snapshots TRUE;mustRecompute
FALSE;lastEngine "MuPAD";xmin "-10";xmax "10";xviewmin "-10";xviewmax
"10";yviewmin "-5";yviewmax "5";viewset"XY";rangeset"X";plottype
4;labeloverrides 1;x-label "t";axesFont "Times New
Roman,12,0000000000,useDefault,normal";numpoints 100;plotstyle
"patch";axesstyle "normal";axestips FALSE;xis \TEXUX{x};var1name
\TEXUX{$x$};function \TEXUX{$\cos \left( 10.0x\right) +\cos \left(
11.x\right) $};linecolor "green";linestyle 1;pointstyle
"point";linethickness 1;lineAttributes "Solid";var1range
"-10,10";num-x-gridlines 900;curveColor "[flat::RGB:0x00008000]";curveStyle
"Line";rangeset"X";function \TEXUX{$2\cos (\frac{10-11}{2}x)$};linecolor
"blue";linestyle 1;pointstyle "point";linethickness 3;lineAttributes
"Solid";var1range "-10,10";num-x-gridlines 100;curveColor
"[flat::RGB:0x000000ff]";curveStyle "Line";VCamFile
'TM4YJ709.xvz';valid_file "T";tempfilename
'TM4YJ700.wmf';tempfile-properties "XPR";}}

\end{document}
